% Options for packages loaded elsewhere
\PassOptionsToPackage{unicode}{hyperref}
\PassOptionsToPackage{hyphens}{url}
\documentclass[
  11pt,
]{article}
\usepackage{xcolor}
\usepackage[margin=1in]{geometry}
\usepackage{amsmath,amssymb}
\setcounter{secnumdepth}{5}
\usepackage{iftex}
\ifPDFTeX
  \usepackage[T1]{fontenc}
  \usepackage[utf8]{inputenc}
  \usepackage{textcomp} % provide euro and other symbols
\else % if luatex or xetex
  \usepackage{unicode-math} % this also loads fontspec
  \defaultfontfeatures{Scale=MatchLowercase}
  \defaultfontfeatures[\rmfamily]{Ligatures=TeX,Scale=1}
\fi
\usepackage{lmodern}
\ifPDFTeX\else
  % xetex/luatex font selection
\fi
% Use upquote if available, for straight quotes in verbatim environments
\IfFileExists{upquote.sty}{\usepackage{upquote}}{}
\IfFileExists{microtype.sty}{% use microtype if available
  \usepackage[]{microtype}
  \UseMicrotypeSet[protrusion]{basicmath} % disable protrusion for tt fonts
}{}
\makeatletter
\@ifundefined{KOMAClassName}{% if non-KOMA class
  \IfFileExists{parskip.sty}{%
    \usepackage{parskip}
  }{% else
    \setlength{\parindent}{0pt}
    \setlength{\parskip}{6pt plus 2pt minus 1pt}}
}{% if KOMA class
  \KOMAoptions{parskip=half}}
\makeatother
\usepackage{longtable,booktabs,array}
\newcounter{none} % for unnumbered tables
\usepackage{calc} % for calculating minipage widths
% Correct order of tables after \paragraph or \subparagraph
\usepackage{etoolbox}
\makeatletter
\patchcmd\longtable{\par}{\if@noskipsec\mbox{}\fi\par}{}{}
\makeatother
% Allow footnotes in longtable head/foot
\IfFileExists{footnotehyper.sty}{\usepackage{footnotehyper}}{\usepackage{footnote}}
\makesavenoteenv{longtable}
\usepackage{graphicx}
\makeatletter
\newsavebox\pandoc@box
\newcommand*\pandocbounded[1]{% scales image to fit in text height/width
  \sbox\pandoc@box{#1}%
  \Gscale@div\@tempa{\textheight}{\dimexpr\ht\pandoc@box+\dp\pandoc@box\relax}%
  \Gscale@div\@tempb{\linewidth}{\wd\pandoc@box}%
  \ifdim\@tempb\p@<\@tempa\p@\let\@tempa\@tempb\fi% select the smaller of both
  \ifdim\@tempa\p@<\p@\scalebox{\@tempa}{\usebox\pandoc@box}%
  \else\usebox{\pandoc@box}%
  \fi%
}
% Set default figure placement to htbp
\def\fps@figure{htbp}
\makeatother
\setlength{\emergencystretch}{3em} % prevent overfull lines
\providecommand{\tightlist}{%
  \setlength{\itemsep}{0pt}\setlength{\parskip}{0pt}}
\usepackage{booktabs}
\usepackage{array}
\usepackage{longtable}
\usepackage{caption}
\usepackage{bookmark}
\IfFileExists{xurl.sty}{\usepackage{xurl}}{} % add URL line breaks if available
\urlstyle{same}
\hypersetup{
  pdftitle={Methods and Computing},
  hidelinks,
  pdfcreator={LaTeX via pandoc}}

\title{Methods and Computing}
\usepackage{etoolbox}
\makeatletter
\providecommand{\subtitle}[1]{% add subtitle to \maketitle
  \apptocmd{\@title}{\par {\large #1 \par}}{}{}
}
\makeatother
\subtitle{DoSS Summer Prep Bootcamp 2026}
\author{}
\date{\vspace{-2.5em}}

\begin{document}
\maketitle

\section{Time and Place}\label{time-and-place}

Exact times TBA

\section{Instructor}\label{instructor}

Benjamin Smith {[}\url{benyamin.smith@mail.utoronto.ca}{]}

\section{Course outline}\label{course-outline}

This course will revisit foundational programming principles in
\texttt{R} and review essential concepts in likelihood inference.

\section{Textbooks}\label{textbooks}

\subsection{Primary Textbooks}\label{primary-textbooks}

\begin{itemize}
\tightlist
\item
  \emph{All of Statistics} by L. Wasserman (AoS)
\item
  \emph{Statistical Inference, Second Edition} by George Casella and
  Robert L. Berger (C\&B)
\end{itemize}

\subsection{Optional Texts}\label{optional-texts}

\begin{itemize}
\tightlist
\item
  \emph{Statistical Models} by A.C. Davison
\item
  \emph{Mathematical Statistics} by K. Knight
\item
  \emph{Theory of Point Estimation} by E.L. Lehmann
\end{itemize}

\section{Tentative Lecture Schedule}\label{tentative-lecture-schedule}

The Lecture Topics and corresponding texts are outlined below. This
schedule is tentative and may be changed as the course progresses.

\newpage

{\def\LTcaptype{none} % do not increment counter
\begin{longtable}[]{@{}
  >{\raggedright\arraybackslash}p{(\linewidth - 4\tabcolsep) * \real{0.2000}}
  >{\raggedright\arraybackslash}p{(\linewidth - 4\tabcolsep) * \real{0.4000}}
  >{\raggedright\arraybackslash}p{(\linewidth - 4\tabcolsep) * \real{0.4000}}@{}}
\toprule\noalign{}
\begin{minipage}[b]{\linewidth}\raggedright
Module
\end{minipage} & \begin{minipage}[b]{\linewidth}\raggedright
Topics
\end{minipage} & \begin{minipage}[b]{\linewidth}\raggedright
References
\end{minipage} \\
\midrule\noalign{}
\endhead
\bottomrule\noalign{}
\endlastfoot
1 & R, Rstudio, and Rmarkdown; Git; data wrangling and graphing; LaTeX;
Elementary data analysis & - \\
2 & Statistical inference (I) & AoS Chp 1-5; AoS Chp 6; C\&B Chp 6.3,
7 \\
3 & Statistical inference (II) & AoS Chp 8; C\&B Chp 8 \\
4 & Linear Regression \& Generalized Linear Models & AoS Chp 13; C\&B
Chp 11-12 \\
5 & Simulation and parallel computing; Bootstrap & C\&B Chap 10; AoS Chp
5, 24; \\
\end{longtable}
}

\end{document}
